% ID: SEL-D0070-OB-E02
% ESTADO: validada. S070.
\textbf{a)} Las dos ramas son suaves en sus respectivos interiores. En $x=1$, $f(1)=3$ y el límite derecho es $1-6+8=3$. El límite izquierdo también vale $3$, así que $f$ es continua en todo $\mathbb R$. Las pendientes laterales en $1$ valen $3\ln3$ y $2-6=-4$. Como difieren, $f$ es derivable salvo en $1$.

\textbf{b)} En $x=3$ se usa la rama cuadrática. Se obtiene $f(3)=9-18+8=-1$ y $f'(3)=6-6=0$. La tangente es horizontal, $\boxed{y=-1}$.
